\section{Síntesis y ampliación}

\subsection{Repaso de los puntos clave}

Esta presentación sistematiza la evolución y el panorama actual de los métodos generativos en 3D:

\begin{enumerate}
    \item \textbf{Evolución de los métodos de generación 3D}: Desde las GANs 3D que requieren ground truth 3D, hasta PlatonicGAN, que utiliza renderizado diferenciable para eliminar la dependencia de datos 3D, y finalmente PixelNeRF, que fusiona múltiples perspectivas.
    \item \textbf{Comprensión 3D en modelos de video}: Los experimentos demuestran que los modelos de video pueden modelar implícitamente la geometría 3D, aunque el costo computacional limita su despliegue.
    \item \textbf{CAT3D}: Paradigma de «difusión multi-visualización + destilación NeRF», capaz de generar escenas 3D en aproximadamente un minuto.
    \item \textbf{B3D}: Optimización mediante una cabeza Gaussian Feedforward que reemplaza a NeRF, generando escenas 3D en menos de 7 segundos.
    \item \textbf{VAE geométrico}: Los VAEs para imágenes no son adecuados para datos geométricos; requieren arquitecturas y funciones de pérdida diseñadas específicamente.
    \item \textbf{Alineación perceptiva VLM}: Utilizar un VLM como juez para entrenamiento DPO, mejorando la calidad perceptiva del autoencoder.
    \item \textbf{ROAR}: Incorporar el relighting como una dimensión de control adicional en los marcos de generación 3D.
    \item \textbf{Captive}: Generación panorámica eficiente en datos mediante cubos de texturas + modelos base congelados.
\end{enumerate}

\subsection{Principios de diseño clave}

\begin{importantbox}{Principios de diseño para métodos de generación 3D}
\textbf{1. Generación 2D + destilación 3D}: Aprovechar las potentes capacidades generativas de los modelos 2D y obtener representaciones 3D eficientes y renderizables mediante reconstrucción 3D.

\textbf{2. Diseño por fases}: Generar propiedades fáciles de supervisar (como RGB, XYZ) con modelos de difusión, y regresionar propiedades difíciles de verificar (como opacidad, escala) con redes feedforward.

\textbf{3. Congelar + adaptar}: Congelar los poderosos sesgos aprendidos en el modelo base y entrenar solo capas de adaptación limitadas para aprender nuevas restricciones.

\textbf{4. Atención multi-visualización}: Garantizar la consistencia 3D en la generación multi-visualización mediante mecanismos de atención global.
\end{importantbox}

\subsection{Lecturas adicionales}

\begin{itemize}
    \item CAT3D: Create Anything in 3D with Multi-View Diffusion Models (Gao et al., Google, 2024)
    \item B3D: Método de generación rápida de escenas 3D feedforward (Google Research)
    \item Splatter Image: Reconstrucción 3D ultra-rápida desde una sola vista (Szymanowicz et al., 2024)
    \item MASt3R: Pipeline de Structure-from-Motion 3D
    \item ROAR: Método de relighting 3D (Google Research)
    \item The Bitter Lesson --- Richard Sutton, 2019
    \item Genie 2: Modelo generativo de video interactivo (Google DeepMind)
    \item World Labs: Modelos mundiales 3D en tiempo real
\end{itemize}

%% --- Fin del contenido --- %%

\end{document}
